\documentclass[12pt]{article}
\usepackage[T1]{fontenc}
\usepackage[utf8]{inputenc}
\usepackage{lmodern}
\usepackage{textcomp}
\usepackage{enumitem}
\usepackage{geometry}
\usepackage{graphicx}
\usepackage{amsmath, amssymb}
\geometry{margin=1in}
\begin{document}
\begin{center}
Political Science - Graduate Level Assessment 23 \
Total Marks: 40 \
Answer all questions.
\end{center}

\section*{Multiple Choice (10 marks)}
\begin{enumerate}[label=\arabic*.]
\item Which of the following might be considered part of the 'deepening and broadening' of security?
\begin{enumerate}[label=(\alph*)]
    \item The nexus between identity and security.
    \item War, the military, and the sovereign state.
    \item Environmental changes.
    \item All of these options.
\end{enumerate}
\item Which of the following concepts does NOT form part of the technical discourse on cyber- security?
\begin{enumerate}[label=(\alph*)]
    \item Business networks
    \item Computer experts
    \item Anti-virus industry
    \item Computer networks
\end{enumerate}
\item Who said "Globalization is not something we can hold off or turn off. It is the economic equivalent of a force of nature-like wind or water"?
\begin{enumerate}[label=(\alph*)]
    \item Ronald Reagan
    \item George Soros
    \item Bill Clinton
    \item George W. Bush
\end{enumerate}
\item What is the referent object in contemporary cyber-security?
\begin{enumerate}[label=(\alph*)]
    \item Digitalised sensitive information
    \item Critical Information Infrastructures
    \item Government IT systems
    \item Telecommunication networks
\end{enumerate}
\item Intergovernmental organizations SELDOM do which of the following?
\begin{enumerate}[label=(\alph*)]
    \item Provide strong enforcement
    \item Monitor parties
    \item Provide fora for discussion
    \item Reduce transaction costs for agreements
\end{enumerate}
\end{enumerate}

\section*{True / False (10 marks)}
\textit{Answer True or False}
\begin{enumerate}[label=\arabic*.]
\setcounter{enumi}{5}
\item True or False: Constraints on US foreign policy decision making include the foreign policies of other states, international law, and intergovernmental organizations.
\item True or False: Rosenau claims that only the President's individual personality is significant in explaining U.S. foreign policy behavior.
\item True or False: The Roosevelt Doctrine emphasized the importance of European intervention in Latin America to maintain stability.
\item True or False: Exceptionalism links isolationist and internationalist strategies as both can be viewed as different means of achieving similar liberal outcomes.
\item True or False: The costs of free trade are often concentrated among specific industries, while the benefits are widely dispersed across the economy.
\end{enumerate}

\section*{Long Form Response (20 marks)}
\textit{Answer each question with a clear and complete explanation.}
\begin{enumerate}[label=\arabic*.]
\setcounter{enumi}{10}
\item Discuss the common criticisms of the human security concept, analyzing their implications for its effectiveness in addressing contemporary global issues.
\item Discuss the interplay between culture and social constructivism in shaping state identity and security perceptions within the realm of political science.
\end{enumerate}

\vfill
\noindent\textit{End of Paper}
\end{document}